\documentclass[11pt]{article}
\usepackage[a4paper,margin=26mm]{geometry}
\usepackage[T1]{fontenc}
\usepackage{lmodern,microtype}
\usepackage{amsmath,amssymb,amsthm,booktabs,array,listings}
\usepackage[hidelinks]{hyperref}
\newtheorem{theorem}{Theorem}[section]
\newtheorem{lemma}[theorem]{Lemma}
\newtheorem{corollary}[theorem]{Corollary}
\newtheorem{proposition}[theorem]{Proposition}
\newcommand{\E}{\mathbb E}
\newcommand{\Prob}{\mathbb P}
\newcommand{\pos}[1]{\left(#1\right)_{+}}
\newcommand{\elln}{\ell_n}
\newcommand{\floor}[1]{\left\lfloor #1\right\rfloor}
\newcommand{\ceil}[1]{\left\lceil #1\right\rceil}
\lstset{basicstyle=\small\ttfamily,columns=fullflexible,keepspaces=true,
  frame=tb,mathescape=true,showstringspaces=false,xleftmargin=1em,
  morekeywords={procedure,while,if,then,else,for,to,do,return,break,end}}
\title{Geometric Target Repair with Biased Sampling\\
\large Logarithmic Height, Increasing Inputs, and Random Permutations}
\author{Vojtech Vincibr}
\date{November 2026}

\begin{document}
\maketitle
\begin{abstract}
We analyze insertion into a binary search tree using a geometrically sampled
\emph{absolute target depth} and a local repair of at most two rotations.
For every fixed $p\in((\sqrt5-1)/2,1)$, every insertion order fixed independently of
the algorithm's randomness gives logarithmic height both in expectation and
with high probability. No balance fields are stored in nodes. For increasing
inputs we give an exact right-spine recurrence, a joint description of the
hanging subtrees, and a deterministic bound on the maximum key's depth.
For uniform random permutations, a separate external-gap potential proves
constant expected ancestor-walk work and exponentially decreasing tails for
rotation distance. These results support a qualified comparison with Brodal's
seven design criteria: the method uses global size information, and constant
structural repair does not by itself make the current Python target calculation
constant time. The repository implements the exact rational choice $p=3/4$.
\end{abstract}

\section{Scope and the seven design criteria}
Brodal~\cite[Introduction]{brodal} studies seven requirements, paraphrased in
Table~\ref{tab:criteria}. His geometric stopping rule measures distance from the
new leaf. Here a geometric draw selects an absolute depth using the current
size. The rotation patterns can agree without the target distributions agreeing.
This paper analyzes the latter algorithm; it does not transfer the former
algorithm's theorems.

\begin{table}[htbp]
\centering\small
\renewcommand{\arraystretch}{1.16}
\begin{tabular}{@{}p{.04\linewidth}>{\raggedright\arraybackslash}p{.31\linewidth}>{\raggedright\arraybackslash}p{.58\linewidth}@{}}
\toprule
 & Criterion (paraphrased) & Result for geometric target repair \\
\midrule
1 & No extra balance information & No per-node fields; a global size counter is required. Thus the strict ``tree and keys only'' condition needs qualification. \\
2 & Constant worst-case rotations & At most two on every input; at most one on monotone inputs. \\
3 & Rotations close to the insertion & Exponential distance tails on monotone inputs and uniform random permutations; no uniform bound conditional on every reachable history. \\
4 & Only local insertion information & Not met: the target uses $n$ and the insertion depth. \\
5 & Constant expected repair time & Proved for structural repair on monotone and uniform random inputs. Full cost also needs the target arithmetic in Section~\ref{sec:arithmetic}. General fixed-order constant repair remains unproved here. \\
6 & Constant expected random bits & At most eight fair bits for either $p=2/3$ or $p=3/4$. \\
7 & Logarithmic expected node depth & A stronger logarithmic expected-height and high-probability height bound holds for every fixed insertion order. \\
\bottomrule
\end{tabular}
\caption{The exact scope of the comparison. No unqualified seven-property claim is made.}
\label{tab:criteria}
\end{table}

\begin{samepage}
Under the relaxed interpretation of item~1 as \emph{no per-node balance fields},
and the arithmetic model of Section~\ref{sec:arithmetic}, with $p=3/4$
(or $p=2/3$), items 1, 2, 3, 5, 6,
and 7 hold for increasing inputs and for uniform random permutations.
This is a six-of-seven statement only with these explicit qualifications.
The general fixed-order height result does not supply the missing general
fixed-order constant-repair result.
\end{samepage}

\section{Model, algorithm, and target law}\label{sec:model}
Start from an empty BST with distinct, totally ordered keys. We may embed their
ranks in the reals so that every gap admits a query value. A successful
insertion first performs ordinary BST leaf insertion, records the new leaf $v$
and its depth $d$, and increments the size to $n$. The root has depth zero;
the height is the largest stored-node depth. Duplicate attempts return before
sampling or changing the size. Nodes contain a key and parent/child links.

Fix $0<p<1$ independently of $n$, and put
\[
  \ell_n=\ceil{\log_{1/p} n}.
\]
Start at $T_n=\ell_n$; repeatedly decrease the target with probability $p$
until the first failure or target zero. All trials are independent, including
across insertions. The mathematical probabilities concern this ideal source;
a finite seeded PRNG provides reproducible executions rather than a proof of
independence. The height theorems assume that the insertion sequence is fixed
independently of these trials. Random permutations below are also independent
of the trial source.

\begin{lstlisting}[caption={Repair after ordinary leaf insertion.}]
procedure Repair($v,n,d,p$)
    $t\gets\lceil\log_{1/p} n\rceil$
    while $t>0$ do
        if Bernoulli($p$) = 0 then break
        $t\gets t-1$
    end while
    if $d\le t$ then return
    if $d=t+1$ then
        RotateUp($v$)
        return
    end if
    $x\gets v$
    for $i=1$ to $d-t-2$ do $x\gets x.\mathrm{parent}$
    $y\gets x.\mathrm{parent}$; $g\gets y.\mathrm{parent}$
    if $(x=y.\mathrm{left})\Longleftrightarrow(y=g.\mathrm{left})$ then
        RotateUp($y$)
    else
        RotateUp($x$); RotateUp($x$)
    end if
end procedure
\end{lstlisting}

\begin{theorem}[Deterministic correctness]\label{thm:correct}
Each successful insertion preserves BST order, stores exactly the old keys and
the new key, and performs at most two rotations. No node balance field is
needed. These statements hold for every sequence of trial outcomes.
\end{theorem}
\begin{proof}
Leaf insertion places the new key at its unique search position. A rotation
preserves the inorder sequence of the rotated subtree and therefore the order
of the entire tree. It changes only structural pointers, including the root
pointer when appropriate. If the two-edge case is reached, $x,y,g$ have depths
$t+2,t+1,t$, so all three exist. On a straight path one rotation promotes $y$;
on a bent path the first rotation makes $x$ a child of $g$, and the second
promotes it above $g$. The remaining cases perform zero or one rotation.
For $n=1$, the target and insertion depth are zero, and no bit is drawn.
\end{proof}

\begin{lemma}[Distribution and sampling cost]\label{lem:target}
Conditional on every history preceding insertion $n$,
\begin{align}
 \Prob(T_n=0)&=p^{\ell_n}, &
 \Prob(T_n=t)&=(1-p)p^{\ell_n-t} \quad(1\le t\le\ell_n),\label{eq:mass}\\
 \Prob(T_n\le t)&=p^{\ell_n-t} \quad(0\le t\le\ell_n), &
 p^{\ell_n}&\le 1/n.\label{eq:cdf}
\end{align}
Writing $G_n=\ell_n-T_n$, we have
\[
 \E G_n=\frac{p(1-p^{\ell_n})}{1-p},\qquad
 \E[\text{number of trials}]=\frac{1-p^{\ell_n}}{1-p}.
\]
\end{lemma}
\begin{proof}
Reaching target $t$ requires $\ell_n-t$ successes. A positive target also
requires a failure; zero stops automatically. This proves the mass and
cumulative formulas. The definition of the ceiling gives $p^{\ell_n}\le1/n$.
The $j$th trial is requested with probability $p^{j-1}$, and
$\Prob(G_n\ge j)=p^j$ for $1\le j\le\ell_n$. Sum the two geometric series.
\end{proof}

For $p=3/4$, read two independent fair bits and succeed unless both are zero.
Every trial reads both bits, so the expected total is
$8(1-(3/4)^{\ell_n})<8$. For $p=2/3$, return failure on \texttt{00}, success
on \texttt{01} or \texttt{10}, and retry on \texttt{11}. The expected bits per
trial are $2/(3/4)=8/3$; the expected total is again less than eight. For general
$p$, the theory assumes exact trials; specifying an arbitrary real probability
on a finite machine requires a representation and sampling model.

\section{Height for every fixed insertion order}\label{sec:height}
Fix a query $q$ never inserted. Its unsuccessful search ends at an external
node (a null child pointer); let $D_n(q)$ be that node's depth, with $D_0(q)=0$.
For a nonempty tree this equals the number of search comparisons. Following
external nodes avoids identifying a newly promoted internal node with the
newly inserted key.

\begin{lemma}[One insertion's local effects]\label{lem:effects}
Write $D,D'$ for the query depths before and after an insertion, including
repair. If the insertion and query use different gaps and diverge at a node
of depth $s$, then for target $t$,
\begin{equation}
 D'-D=\begin{cases}
  1 & t=s,\\
  0 & t=s-1,\\
 -1 & 0\le t\le s-2,\\
  0 & t>s.
 \end{cases}                                                   \label{eq:effects}
\end{equation}
If the insertion is into the query's own gap, then $D'\le D+1$.
If in addition $D\ge\ell_n+2$, then $D'=D$ for every target.
\end{lemma}
\begin{proof}
For different gaps the leaf insertion itself leaves $D$ unchanged. Consider
$g,y,x$ at depths $t,t+1,t+2$. A query that diverges at $g$ lies in its child
subtree opposite $y$, which moves down one level in either repair. Divergence
at $y$ places the query in the subtree opposite $x$, whose depth is unchanged.
Agreement through $x$ places the query in an original child subtree of $x$;
either repair moves that subtree up one level. In a bent repair $x$ moves up
two levels but its original child subtrees move up only one. A divergence
above $g$ is unaffected. If repair is just a single promotion of the new leaf,
the opposite subtree of its parent moves down one; earlier divergences are
unaffected. If repair does nothing, then $t\ge d>s$.

For the same gap, the inserted leaf replaces one external node of depth $D$
by two of depth $D+1$. Neither repair pushes either new external node farther
down. If $D\ge\ell_n+2$, every target is at most $D-2$. When $D>t+2$,
both new external nodes lie below $x$ and move up one level. When $D=t+2$,
$x$ is the new leaf: a straight repair again moves both up one; a bent repair
makes them the inner null pointers of the two old ancestors at depth $D$.
\end{proof}

For the rest of the probabilistic depth analysis assume
\begin{equation}
 p>\frac{\sqrt5-1}{2},\qquad
 a=\frac{p^2}{1-p}>1,\qquad
 \alpha=\log_{1/p}a,\qquad \beta=1+\alpha.
                                                               \label{eq:constants}
\end{equation}
Constants in asymptotic notation may depend on this fixed $p$.

\begin{lemma}[Conditional exponential drift]\label{lem:drift}
For a fixed input sequence and its history $\mathcal F_{n-1}$,
\begin{equation}
 \E[a^{D_n(q)}\mid\mathcal F_{n-1}]
 \le \left(1+\frac{a-1}{n}\right)a^{D_{n-1}(q)}+a^{\ell_n+2}.
                                                               \label{eq:drift}
\end{equation}
For a different-gap insertion, the additive term can be omitted.
\end{lemma}
\begin{proof}
Condition on the history, so the tree and next key are fixed. In a different
gap with $2\le s\le\ell_n$, let $A$ and $B$ be the probabilities that the query
depth increases and decreases. Lemmas~\ref{lem:target} and~\ref{lem:effects} give
\[
 A=(1-p)p^{\ell_n-s},\qquad B=p^{\ell_n-s+2}=aA.
\]
Consequently
\begin{equation}
 \E[a^{D'-D}]=1+A(a-1)+B(a^{-1}-1)=1.                         \label{eq:cancel}
\end{equation}
If $s>\ell_n$, depth cannot increase. At $s=0$ its increase probability is
$p^{\ell_n}\le1/n$. At $s=1$ it is either zero or
$((1-p)/p)p^{\ell_n}\le1/n$, since $p>1/2$. These cases give the
multiplicative factor in~\eqref{eq:drift}. For a same-gap insertion,
$D\ge\ell_n+2$ gives $D'=D$; otherwise $D'\le\ell_n+2$. The nonnegative
additive allowance covers the latter case.
\end{proof}

\begin{lemma}[Polynomial moment]\label{lem:moment}
For every fixed query $q$ outside the input and every $n\ge1$,
\begin{equation}
 \E[a^{D_n(q)}]\le a^3 n^\beta.                              \label{eq:moment}
\end{equation}
\end{lemma}
\begin{proof}
First $\beta>a$. Indeed,
\[
 \beta-1=\frac{\ln a}{-\ln p}
 >\frac{a-1}{a}\frac{p}{1-p}=\frac{a-1}{p}>a-1,
\]
using $\ln a\ge(a-1)/a$ and $-\ln p<(1-p)/p$.
The ceiling implies $a^{\ell_n+2}\le a^3n^{\beta-1}$. Taking expectations
in~\eqref{eq:drift} gives, with $M_n=\E[a^{D_n(q)}]$,
\[
 M_n\le(1+(a-1)/n)M_{n-1}+a^3n^{\beta-1}.
\]
For $n\ge2$, logarithms and $\ln(1+u)\le u$, $\ln(1-u)\le-u$ show
\[
 (1+(a-1)/n)(1-1/n)^{\beta-1}\le1.
\]
Multiplying by $1-1/n$ yields
\[
 (1+(a-1)/n)a^3(n-1)^\beta
 \le a^3n^\beta-a^3n^{\beta-1}.
\]
This absorbs the additive term in the induction. The base is $M_1=a\le a^3$.
\end{proof}

\begin{theorem}[Height, expectation, and high probability]\label{thm:height}
After any fixed sequence of $n\ge1$ distinct keys, the height $H_n$ satisfies
\begin{equation}
 \Prob(H_n\ge h)\le
 \min\{1,\,2a^2n^{\beta+1}a^{-h}\}\qquad(h\in\mathbb Z_{\ge0}).
                                                               \label{eq:height-tail}
\end{equation}
In particular, $\E H_n=\Theta(\log(n+1))$ as $n\to\infty$.
For every $\gamma>0$ there is a constant $c_{p,\gamma}$ such that
$\Prob(H_n>c_{p,\gamma}\ln n)\le n^{-\gamma}$ for $n\ge2$.
Every stored node has expected depth $O(\log(n+1))$.
\end{theorem}
\begin{proof}
Choose one fixed query in each of the $n+1$ gaps of the final key set, before
any algorithmic randomness. They reach every final external node, so
$H_n+1=\max_q D_n(q)$. Markov's inequality and a union bound give
\[
 \Prob(H_n\ge h)
 \le(n+1)a^3n^\beta a^{-(h+1)}
 \le2a^2n^{\beta+1}a^{-h}.
\]
There is no independence assumption between queries. Taking
$h\ge(\beta+1+\gamma)\log_a n+\log_a(2a^2)$ proves the high-probability claim.
For $h_0=\ceil{\log_a(2a^2n^{\beta+1})}$, the tail-sum formula gives
\[
 \E H_n\le h_0+\sum_{j\ge1}a^{-j}=h_0+\frac1{a-1}=O(\log(n+1)).
\]
Every binary tree satisfies $n\le2^{H_n+1}-1$, which gives the asymptotic
matching lower bound. Stored-node depths are bounded by height pointwise.
\end{proof}

\section{Increasing sequences: exact dynamics}\label{sec:increasing}
The \emph{right spine} follows right children from the root to the current
maximum. Let $R_n$ be its number of vertices after $n$ insertions, with $R_0=0$.
All search edges for an increasing insertion go right. Thus any two-edge repair
is straight; a repair moves the new maximum up exactly one level, even if the
promoted node is an older ancestor.

\begin{theorem}[Spine and local cost on monotone inputs]\label{thm:monotone}
For increasing inputs, at every insertion
\begin{equation}
 R_n=R_{n-1}+\mathbf1_{\{T_n\ge R_{n-1}\}},\qquad
 R_n-1\le\ell_n.                                             \label{eq:spine}
\end{equation}
At most one rotation is used. If $V_n$ is the implementation's number of counted
ancestor visits (the walked departures plus the final inspected ancestor), then
\begin{equation}
 V_n=(d-T_n-1)_+,\quad V_n\le G_n,\quad
 \E[V_n\mid\text{past}]\le\frac{p(1-p^{\ell_n})}{1-p}.
                                                               \label{eq:mono-visits}
\end{equation}
For every integer $k\ge1$, $\Prob(V_n\ge k\mid\text{past})\le p^k$.
The same claims hold with left and right exchanged for decreasing inputs.
\end{theorem}
\begin{proof}
The new maximum attaches at depth $d=R_{n-1}$. If $T_n\ge d$ no rotation
occurs, and the spine gains one vertex. Otherwise one rotation removes exactly
one old spine vertex from the extended spine, giving~\eqref{eq:spine}.
For $n=1$ the spine has one vertex at depth zero. For $n\ge2$,
inductively $d\le\ell_{n-1}+1\le\ell_n+1$. With no repair $d\le T_n\le\ell_n$;
with repair the new maximum's final depth is $d-1\le\ell_n$.
For arbitrary inputs the parent loop has $d-T_n-2$ steps if that quantity is
nonnegative and the two-edge case is reached, followed by one counted visit.
Hence $V_n=(d-T_n-1)_+$. The bound on $d$ gives $V_n\le\ell_n-T_n=G_n$.
Lemma~\ref{lem:target} supplies its mean and tail.
\end{proof}

Set $\pi_n(j)=\Prob(R_n=j)$ and extend it by zero outside its support. For integer $u$, define
\[
 F_n(u)=\Prob(T_n\le u)=
 \begin{cases}0&u<0,\\p^{\ell_n-u}&0\le u\le\ell_n,\\1&u>\ell_n.\end{cases}
\]
Independence of the fresh target gives the exact recurrence
\begin{equation}
 \boxed{\ \pi_n(j)=\pi_{n-1}(j)F_n(j-1)
       +\pi_{n-1}(j-1)\bigl(1-F_n(j-2)\bigr)\ },             \label{eq:spine-law}
\end{equation}
with $\pi_0(0)=1$. The expected maximum-key depth is
$\sum_j(j-1)\pi_n(j)$. This recurrence keeps or increases spine length by one;
a single rotation does \emph{not} reset it to the target depth.
If $A_n$ is the total number of rotations through insertion $n$, then summing
one added spine vertex per insertion and one removed vertex per rotation gives
$A_n=n-R_n$. Hence
\[
 \max\{0,n-\ell_n-1\}\le A_n\le n-1,
 \qquad A_n/n=1-O_p(\log(n+1)/n).
\]
Thus average rotations approach one deterministically on increasing inputs.
The current maximum's depth bound does not bound the whole tree deterministically:
all-zero targets can produce a left chain of height $n-1$.

The hanging left subtrees require a joint state. Write the spine as
$a_0,\ldots,a_{r-1}$ and its left subtrees as $S_0,\ldots,S_{r-1}$, and append
the new maximum $a_r$ with $S_r=\varnothing$. For $t<r$, remove $a_t$ from
the spine and replace the promoted node $a_{t+1}$'s left subtree by
\begin{equation}
 \operatorname{Node}(a_t,S_t,S_{t+1}).                         \label{eq:merge}
\end{equation}
Later spine entries shift up one position. If $t\ge r$ there is no merge.
With empty-subtree height $-1$, the merged height is
$1+\max\{\operatorname{height}(S_t),\operatorname{height}(S_{t+1})\}$.
The tuple of spine keys and subtrees, together with~\eqref{eq:mass}, is a closed
transition description. Individual subtree-height marginals cannot be multiplied
as if independent.

For example, insert $1,2,3,4$ with $T_2\ge1,T_3=1,T_4=0$. The result has
root $3$, right child $4$, and left subtree with root $1$ and right child $2$.
A hanging left subtree therefore need not be a left-going chain.
Theorem~\ref{thm:height} bounds the height of all these subtrees together,
not only the right spine, by $O(\log n)$ in expectation and with high probability.

\section{Uniform random permutations}\label{sec:random}
Let the input be a uniform permutation of $N$ fixed distinct keys, independent
of the repair trials. All bounds below hold at every prefix $1\le n\le N$.
Conditioning on the entire permutation immediately transfers
Theorem~\ref{thm:height} to this model. To bound \emph{repair work}, we need a
stronger estimate for the next insertion gap.

\begin{lemma}[The next relative rank is uniform]\label{lem:rank}
Conditional on the previous relative insertion ranks and repair trials, the
next key enters each of the $n$ external gaps of the $(n-1)$-node tree with
probability $1/n$.
\end{lemma}
\begin{proof}
Let $J_i$ be the rank of the $i$th key among the first $i$ keys. The map from
permutations to vectors $(J_1,\ldots,J_N)$ with $1\le J_i\le i$ is a bijection:
recover the last key as the $J_N$th remaining key, remove it, and continue
backwards. Thus the $J_i$ are independent and uniform. The tree shape with keys
replaced by their current ranks is determined by earlier $J_i$ and earlier
repair trials. Each possible new rank corresponds to one of its external gaps.
This conditioning is on comparison ranks, \emph{not} on the actual numerical
prefix keys; conditioning on those values need not make the remaining gaps
equally likely.
\end{proof}

\begin{theorem}[External-gap potential and insertion depth]\label{thm:random-potential}
Let $\mathcal G_n$ be all $n+1$ external gaps after insertion $n$, and set
\[
 Q_n=\sum_{g\in\mathcal G_n}a^{\operatorname{depth}(g)},\qquad
 C_p=\max\left\{1,\frac{2a^2(a-1)}{\beta-a}\right\}.
\]
Then $\E Q_n\le C_p(n+1)^\beta$ for all $n\ge0$.
If $I_n$ denotes the new leaf's depth \emph{before} repair, then
\begin{align}
 \E[a^{I_n-\ell_n}]&\le C_p,\label{eq:excess-moment}\\
 \Prob(I_n\ge\ell_n+k)&\le C_p a^{-k}\quad(k\ge0),\label{eq:excess-tail}\\
 \E I_n&\le\log_{1/p}n+\log_a C_p.\label{eq:insert-depth}
\end{align}
\end{theorem}
\begin{proof}
Condition on the ranked tree and on the chosen insertion gap, of depth $d$.
For each other gap, Lemma~\ref{lem:drift} bounds its expected exponential depth
by $(1+(a-1)/n)$ times its old value. The insertion gap is replaced by two gaps.
If $d\ge\ell_n+2$, both new depths are exactly $d$ by Lemma~\ref{lem:effects}.
Otherwise each is at most $d+1$, with $d\le\ell_n+1$. Their combined
contribution is therefore always at most
\[
 2a^d+2(a-1)a^{\ell_n+1}.
\]
Put $c=a-1$. Before averaging over the uniform chosen gap, the expected total
is at most
\[
 (1+c/n)(Q_{n-1}-a^d)+2a^d+2c\,a^{\ell_n+1}.
\]
Lemma~\ref{lem:rank} gives $\E[a^d\mid\text{ranked tree}]=Q_{n-1}/n$, hence
\begin{align*}
 \E Q_n
 &\le(1+a/n-c/n^2)\E Q_{n-1}+2c\,a^{\ell_n+1}\\
 &\le(1+a/n)\E Q_{n-1}+2a^2(a-1)n^{\beta-1}.
\end{align*}
Here $a^{\ell_n+1}\le a^2 n^{\beta-1}$. Since $\beta>a>1$, convexity gives
$(n+1)^\beta\ge n^\beta+\beta n^{\beta-1}$. The induction hypothesis
$\E Q_{n-1}\le C_p n^\beta$ and the definition of $C_p$ thus imply
$\E Q_n\le C_p(n+1)^\beta$. Its base is $Q_0=1$.

The next gap is uniform, so
$\E a^{I_n}=\E Q_{n-1}/n\le C_p n^{\beta-1}$.
Since $a^{\ell_n}\ge n^{\beta-1}$, this proves~\eqref{eq:excess-moment};
Markov proves~\eqref{eq:excess-tail}. Jensen's inequality applied before
rounding the logarithm gives~\eqref{eq:insert-depth}.
\end{proof}

\begin{corollary}[Constant structural repair and rotation locality]\label{cor:random-work}
For uniform random permutations, counted ancestor visits satisfy
\begin{equation}
 \E V_n\le \frac{C_p}{a(a-1)}+\frac{p}{1-p}=O_p(1).           \label{eq:random-visits}
\end{equation}
For every integer $k\ge1$,
\begin{equation}
 \Prob(V_n\ge k)\le C_p a^{-k/2}+p^{k/2}.                    \label{eq:visit-tail}
\end{equation}
The same tail bounds the largest original distance from the new leaf to a node
promoted during that insertion (distance zero when no rotation occurs).
Sampling, pointer walking, and rotations together take expected $O_p(1)$ time,
excluding the separate computation of $\ell_n$ and the insertion search.
\end{corollary}
\begin{proof}
Pointwise, $V_n=(I_n-T_n-1)_+\le(I_n-\ell_n-1)_++G_n$.
By~\eqref{eq:excess-tail} and the tail-sum formula,
\[
 \E(I_n-\ell_n-1)_+
 =\sum_{j\ge1}\Prob(I_n\ge\ell_n+1+j)
 \le C_p\sum_{j\ge1}a^{-1-j}=\frac{C_p}{a(a-1)}.
\]
Lemma~\ref{lem:target} proves~\eqref{eq:random-visits}. If $V_n\ge k$,
then $I_n-\ell_n+G_n\ge k+1$. At least one summand exceeds $k/2$;
Markov for the first and the geometric tail for the second give
\eqref{eq:visit-tail}. No independence between these events is needed.
A straight repair promotes a node at original distance $I_n-T_n-1=V_n$;
a bent repair promotes one at distance $I_n-T_n-2=V_n-1$.
Single-leaf promotion has distance zero. Lemma~\ref{lem:target} and
Theorem~\ref{thm:correct} bound the remaining structural work.
\end{proof}

The potential also gives a sharper random-order height tail:
\begin{equation}
 \Prob(H_n\ge h)\le
 \min\{1,\ C_p(n+1)^\beta a^{-(h+1)}\},                     \label{eq:random-height}
\end{equation}
because an external node of depth at least $h+1$ contributes at least $a^{h+1}$
to $Q_n$. For $p=3/4$ the constants are $a=9/4$,
$\beta=1+\log_{4/3}(9/4)$, and~\eqref{eq:random-visits} is less than six
counted visits. For an exact check of this numerical claim,
$(9/4)^4>(4/3)^{11}$ gives $\beta>15/4$, so $C_p<135/16$ and
$C_p/[a(a-1)]<3$. These are conservative upper bounds, not predicted measurements.

\begin{proposition}[The shape law is not the ordinary random-BST law]\label{prop:shape}
For a uniform permutation of three keys, the probability of the balanced
three-node tree is $(1+p^{\ell_3})/3$. With $p=3/4$, it is $337/768$,
whereas ordinary unbalanced insertion gives probability $1/3$.
\end{proposition}
\begin{proof}
Conditional on the missing third key, the root of the two-key tree is uniform
among those two keys, by symmetry of their insertion orders and of repair.
The last key is the median with probability $1/3$. Its path bends at depth two,
and the result is balanced exactly when $T_3=0$.
If the last key is an extreme, the old root is the final median with probability
$1/2$: then the initial three-node tree is balanced and stays so exactly when
$T_3\ne0$. With the other old root, the path is straight and becomes balanced
exactly when $T_3=0$. The extreme case therefore has balance probability $1/2$.
The total is $p^{\ell_3}/3+(2/3)(1/2)$. Ordinary insertion is balanced exactly
when its first key is the median. Finally $\ell_3=4$ at $p=3/4$.
\end{proof}

\section{Costs, arithmetic, and limits of the claims}\label{sec:arithmetic}
For arbitrary inputs the exact counted visit formula gives
\begin{equation}
 \E(I_n-\ell_n-1)_+\ \le\ \E V_n\ \le\
 \E(I_n-\ell_n-1)_++\frac{p}{1-p}.                           \label{eq:excess}
\end{equation}
Thus constant expected structural repair is equivalent to constant expected
positive excess above the reference depth, up to additive constants. An
$O(\log n)$ height bound alone does not imply this excess bound.
Theorems~\ref{thm:monotone} and~\ref{thm:random-potential} supply it for the two
input classes studied here. For arbitrary fixed inputs we establish expected
$O(\log(n+1))$ search and structural repair, using Theorem~\ref{thm:height}
on the previous prefix. The deterministic two-rotation bound always holds.

\paragraph{Computing the ceiling exactly.}
For the implemented $p=3/4$, the maximum target is the least integer $k$ for
which $4^k\ge n3^k$. Python recomputes these powers in $\ell_n$ loop iterations;
the C++ implementation searches a finite table of exact thresholds
$\floor{(4/3)^k}$ for its 32-bit node-count range. The table covers that finite
range only. Neither implementation justifies an unbounded constant-time
arithmetic claim solely from the structural-repair theorem.

There is a simple exact arithmetic variant for any fixed rational $p=u/v$,
$0<u<v$: retain global values $k,A=v^k,B=u^k$, initially $0,1,1$. After each
size increment, while $A<nB$, replace $(k,A,B)$ by $(k+1,vA,uB)$. The result is
$k=\ell_n$. At insertion $n\ge2$ the number of iterations is at most
$\ceil{\log_{v/u}2}$, because $n/(n-1)\le2$. For fixed $u,v$, $A$ and $B$
have $O(\log n)$ bits, hence occupy a constant number of words of size
$\Theta(\log n)$. In a word-RAM with constant-time word multiplication this
update costs $O_p(1)$ time and adds only a constant number of global words.
It changes no target law or node representation. This is an explicitly stated
arithmetic variant, not a claim that the present Python loop already has that
bound. Bit-complexity costs of arbitrary-precision integers are different.

With that arithmetic variant, full rebalancing after the ordinary search has
supplied both the leaf position and its depth has expected $O_p(1)$ time on
monotone inputs and uniform random permutations. A leaf position alone does not
supply its depth: walking to the root to recover it can take nonconstant time.
Locating that position still takes expected $O(\log(n+1))$ time. For arbitrary
fixed inputs the complete insertion likewise takes expected $O(\log(n+1))$
time in this model.

\paragraph{Conditioning and adaptive inputs.}
There is no uniform constant repair bound after every reachable history.
Insert $1,2,\ldots,m$ and let every target be zero, a positive-probability event.
Inductively the new maximum becomes root, producing the left chain
$m,m-1,\ldots,1$. Inserting $0$ now has depth $m$, so every target gives
$V\ge(m-\ell_{m+1}-1)_+$. This is a conditional lower bound on a rare history,
not a lower bound on the unconditional expectation for that fixed input.

For keys selected after observing the tree, the final gaps used in the union
bound need not be fixed independently of the random choices. The uniform-rank
lemma likewise fails for a general adaptive rule. No adaptive-input height
bound is claimed here. The historical fair-coin paper in this repository
studies $p=1/2$ with a different cap and cannot be used as the new biased
strategy's height proof: at $p=1/2$, $a<1$, so the increasing exponential
potential above is unavailable.

\section{Implementation correspondence and reproducibility}\label{sec:reproduce}
The registered strategy is \texttt{geometric-three-quarters}.
Both the Python and C++ implementations use the law in Lemma~\ref{lem:target},
read exactly two bits per trial, stop at target zero, and call the same local
repair pattern. The first insertion consumes no bits. The fixed $p$ is shared
across insertions. Native claims apply while the node count fits its supported
index range. The browser implementation follows the same target law.

The repository checks identical scripted-bit traces across implementations,
exact integer depth boundaries, and local repair against independent immutable
rotations. Additional bounded tests enumerate external-gap depth changes and
small exact distributions. These tests check correspondence on finite inputs;
the asymptotic statements follow from the mathematical proofs above. The
existing Lean development covers selected deterministic repair and fair-coin
sampler facts; it does not machine-check the new stochastic height or
random-permutation theorems in this paper.

Compile from the repository root with:
\begin{lstlisting}[mathescape=false,basicstyle=\footnotesize\ttfamily]
latexmk -pdf -interaction=nonstopmode -halt-on-error \
  -outdir=/tmp/geometric-biased-repair-paper \
  docs/proofs/geometric-biased-repair.tex
\end{lstlisting}
GitHub Actions compiles this source, publishes a downloadable workflow artifact,
and includes the newly compiled PDF in the GitHub Pages site. The repository's
proof README gives the local validation commands and the two PDF destinations.

\begin{thebibliography}{1}
\bibitem{brodal}
Gerth St\o lting Brodal.
\emph{Bottom-up Rebalancing Binary Search Trees by Flipping a Coin}.
Full version, 2024, arXiv:2404.08287v1.
\url{https://arxiv.org/abs/2404.08287v1}.
\end{thebibliography}
\end{document}
